\documentclass{article}
\usepackage{graphicx}
\usepackage{xcolor}
\usepackage{soul}
\usepackage{newunicodechar}
\usepackage{tikz}
\usetikzlibrary{positioning}
\usepackage{fontspec}
\usepackage{polyglossia}
\setmainlanguage{english}
\setotherlanguage{greek}
\usepackage{geometry}
\geometry{margin=1in}
\usepackage{csquotes}
\usepackage{enumitem}
\usetikzlibrary{arrows.meta, positioning, shapes.geometric}
\usepackage{fancyvrb}
\usepackage{setspace}
\newfontfamily\greekfont{Gentium Plus}
\newcommand{\gk}[1]{{\greekfont #1}}
\newunicodechar{₀}{$_0$}
\newunicodechar{₁}{$_1$}
\newunicodechar{₂}{$_2$}
\newunicodechar{₃}{$_3$}
\newunicodechar{₄}{$_4$}
\usepackage[colorlinks=true, linkcolor=blue, urlcolor=blue, citecolor=blue]{hyperref}
\usepackage{changepage}
\usepackage{comment}
\newcommand{\ra}{\ensuremath{\rightarrow}\ }
\newcommand{\inlinenote}[1]{\par\noindent
  \fcolorbox{orange}{orange!20}{\parbox{0.97\linewidth}{#1}}\par\noindent}

% ============================================================================
% DRAFT-v3 — Part II Methodology, "The Justification of the Method".
% Justification layer COMPLETE: five movements (orientation; the producing/using
%   art; rhetorical incorporation; Calleja; Burke). Ported faithfully from
%   DRAFT-v3.md (per-paragraph-revised to the Part I ¶3 voice; the two poles =
%   world-disclosedness / co-disclosedness, Heidegger's own terms, BT 145+160).
% Secondary lit cited inline (Author, ShortTitle, Page); Greek via \gk{}+Bekker;
%   nodes A_0..A_4 in math mode. Unused biblatex dropped for standalone compile
%   (no \cite in body); re-add when integrating into the dissertation tree.
% DEFERRED: M.5 instruments+dwelling · M.6 VR<->screen differential ·
%   M.7 protocol · Appendices A/B/C.
% ============================================================================

\begin{document}

\section{The Justification of the Method}

\subsection{Rhetorical-Ontological Diachronic Analysis: The Mechanism of a Reshaping}

The preceding account did one thing only. It held that the virtual and the phantasmic are real---distinct in kind both from the non-virtual and from the grounds on which they hang---and that they share one mode of being: a potency rendered into act. That much is settled. What it set aside---and set aside deliberately---was a question of quite another order: not what the rendered world \textit{is}, but how it comes to work upon, and in time to remake, the one who enters it. The first question was metaphysical. The second is rhetorical---and it calls for a method, not a further ontology. That method is the one I develop here, and I call it Rhetorical-Ontological Diachronic Analysis (RODA).

What RODA does is best said through what it reads. It takes a stretch of play---a bounded run of engagement, motion by motion---and reads it as the \textit{actualization of desire}: the passage of a desiring being from the first stir of perception, through a charged image and a committed judgment, to a completed act. That passage is the object the method traces. It then asks how a designed world drives that passage toward a substantial union with itself---the union I call \textit{rhetorical incorporation}. And it asks what that union becomes once it is won: it does not lapse when the screen goes dark, but settles into the player's \textit{hexis} (\gk{ἕξις}, \textit{NE} II.1 1103a16), the standing disposition from which the next encounter sets out. The wager is that engaging and reshaping are not two events but one: the second is the sediment of the first, so that to trace the actualization is already to trace the remaking.

Three words in the method's name mark the three things it must do, and each governs a movement below. The analysis is \textit{diachronic}: it does not photograph a state of involvement but follows a motion through time. It is read forward, then---motion by motion---even as it is understood backward from its end. It is \textit{ontological} before it is psychological: the union it traces is a convergence in being, not a swell of feeling, and for that it must borrow an ontology rather than a vocabulary of affect. And it is \textit{rhetorical}: a world does not merely stand before the player but moves the player, and the grammar of such moving---how an act is given its motive, and how that motive's true source is concealed---belongs to rhetoric as the ancients conceived it. Where the metaphysics carried the reality-burden, the method carries the mechanism.

\subsection{The Producing Art and the Using Art: What the World Supplies and What the Player Does}

The frame beneath the chain is Part I's, not this section's, and I rehearse only the two features the argument leans on. The first is the unity of the change. It is tempting to think that, because a played world has both a maker and a player, there are two changes, each with its own full complement of causes; but the four causes belong to a single change, and they cannot be parcelled out between two agents. Designing and playing are indeed two acts---yet they are joined. They are joined at the relation Aristotle draws between the producing art and the using art (\textit{Phys.} II.2, 194a34--b5): the art that makes the flute is not the art that plays it, and the maker knows the matter while the user knows the form. From this the partition falls out clean. What the world \textit{supplies}---its matter and its form---is the maker's finished product handed over as a given; what the player \textit{does}---the bringing-about and the for-the-sake-of-which---is the using art worked upon that given. The crossing is exact: the designer's end-product is the player's given means. To say that a world \textit{engages} the player, then, names neither a bare meeting nor a metaphor but this precise handing-over, in which the player's own efficient and final causation takes up the world's material and formal supply as the stuff of its act. The designer cannot supply the player's desire, and so cannot be the efficient cause of the player's act; the designer can only shape the world the player meets. Why the act comes \textit{now} is the designer's doing; \textit{which way} it runs is the player's own.

The second feature is the shape of the union toward which the change tends. Part I took this shape from Aristotle's account of perceiver and perceived: ``the activity of the sensible object and that of the sense is one and the same activity, and yet the distinction between their being remains'' (\textit{DA} III.2, 425b26--27). This is the union-structure the method inherits---one in substrate, different in being---and it must be held with care, for the whole third movement turns on it. It does not say the two collapse into one thing. It says that one actuality is at once the actuality of both, while what it is to be the one differs from what it is to be the other. Agent and world stand so: correlative poles of a single actualization, joined in act, separate in being. This formula, and only this, lets the method speak of a \textit{substantial} union without the absurd consequence that player and world become numerically one.

\subsection{Rhetorical Incorporation: The Substantial Union of Player and World}

The spine of the method is the chain---the actualization-of-desire chain. It is read forward, motion by motion, yet understood backward from the \textit{orekton} (\gk{ὀρεκτόν}), the desired object that draws the whole toward itself and lends each earlier node its sense, as an end lends sense to a means. It opens at $A_0$, the ontological horizon: a sensible object and the sense-faculty given together within motion and time, the designed \textit{Umwelt} into which the player is always already thrown. The perceptual motion $M_0 \rightarrow A_1$ reaches $A_1$, completed perception, which lays down a dual residue---the \textit{resonant kinēsis}, the after-motion that lingers when the object withdraws, in its two strands: the \textit{resonant aisthēma}, the persisting eidetic trace, and the \textit{resonant epithymia}, the persisting hedonic-orectic tone. The phantastic motion $M_1 \rightarrow A_2$ renders, from that residue, $A_2$, the \textit{phantasma} proper---the determinate, available, charged image on which every later operation works, the pivot of the chain. The cognitive motion $M_2 \rightarrow A_3$ carries the image to $A_3$, completed cognition: the \textit{phantasma} taken up by the three orientational modes---intellection, memory, the discursive-deliberative---and ratified by the orthogonal committal \textit{doxa} (\gk{δόξα}), the taking-as-true, in which the emotion concretizes. And the orektikon motion $M_3 \rightarrow A_4$ issues in $A_4$, the completed act, which sediments into \textit{hexis}. That cross-pass settling is \textit{temporal sedimentation}, and it must not be confused with the within-pass residue of $A_1$; the two are different deposits at different scales.

The method's central event falls at the $A_3 \rightarrow A_4$ transition, where committed cognition passes into act, and it is in just this passage that the union is consummated. It has two correlative poles, held not as two events but as one: \textit{world-disclosedness} (\textit{Welterschlossenheit}) and \textit{co-disclosedness} (\textit{Miterschlossenheit}). The parallel in the names is the claim itself---one disclosedness, opened at once toward a world and toward others, as one actuality is at once the actuality of both: one in substrate, different in being. In world-disclosedness the player is taken into the world acted in, and the world into the player's awareness---the designed \textit{Umwelt} disclosed as a place dwelt in, not a scene surveyed. In co-disclosedness the same act is an acting-together, a being-with those the player acts alongside---a ``we'' the act at once presupposes and ratifies. The poles are not successive and not separable: the player drawn into the world is, in the very same motion, drawn into a being-with-others within it.\footnote{Both names are Heidegger's. \textit{Welterschlossenheit} is the disclosedness of world---its worldhood; \textit{Miterschlossenheit}, which Macquarrie and Robinson render ``co-disclosedness,'' he states of space, ``the essential co-disclosedness of space'' (\textit{BT} 145), and I extend it to the co-disclosedness of \textit{others}. \textit{BT} 160 holds the two together: with Dasein's Being-with, the others' ``disclosedness has been constituted beforehand,'' and this ``goes to make up significance---that is to say, worldhood.'' World-disclosedness and co-disclosedness are there co-constituted---disclosed in one stroke, not in succession.}

The convergence is experiential and ontological, an oneness \textit{in act}---and it is emphatically not a change of kind. The player wholly taken into a virtual world has not thereby made that world non-virtual, no more than the perceiver who is one-in-substrate with the perceived has turned the seeing into the thing seen. The reality-burden secured the virtual as a distinct kind of being; rhetorical incorporation does not undo that distinctness but presupposes it.

\subsection{Incorporation, Re-Grounded: Calleja and the Channels of Involvement}

The chain shows that a designed world engages the player's perception, cognition, and desire. What it cannot say, from its own resources, is \textit{where} in the world's design that engaging happens---nor \textit{whether}, in a given case, the player has truly been taken into the world rather than merely set before it. For both the dissertation turns to Gordon Calleja, and it turns to him because his is a vocabulary built from games and tested against players, not deduced from a philosophy of mind. The turning is a sublation, in three beats---a preserving, a transforming, an elevating.

What is \textit{preserved} is incorporation itself, kept as world-disclosedness, together with the six dimensions of Calleja's Player Involvement Model, taken up as so many designed channels through which a world reaches the player along the chain.\footnote{The six dimensions of the Player Involvement Model (Calleja, \textit{In-Game}, 43--44; full verbatim articulation in the Calleja appendix): \textit{kinesthetic}, ``all modes of avatar or game piece control\ldots{} ranging from learning controls to the fluency of internalized movement''; \textit{spatial}, ``engagement with the spatial qualities of a virtual environment\ldots{} the sense that they are inhabiting a place, rather than merely perceiving a representation of space''; \textit{shared}, ``engagement derived from players' awareness of and interaction with other agents\ldots{} cohabitation, cooperation, and competition''; \textit{narrative}, ``the narrative that is scripted into the game and the narrative that is generated from the ongoing interaction with the game world, its embedded objects and inhabitants, and the events that occur there''; \textit{affective}, ``various forms of emotional engagement\ldots{} either purposefully designed into the game or precipitated by an individual player's interpretation of in-game events and interactions with other players''; \textit{ludic}, ``players' engagement with the choices made in the game and the repercussions of those choices.''} Calleja's own definition is kept whole: incorporation is ``the absorption of a virtual environment into consciousness, yielding a sense of habitation, which is supported by the systemically upheld embodiment of the player in a single location, as represented by the avatar'' (Calleja, \textit{In-Game}, 169); and its double axis with it---the player ``incorporates\ldots{} the game environment into consciousness while simultaneously being incorporated through the avatar into that environment,'' where ``the simultaneous occurrence of these two processes is a necessary condition for the experience of incorporation'' (169). The channels name where the world's supply presses on the chain; they are named in an analysis only where the world is in fact pressing, never run down as a list to fill.

What is \textit{transformed} is the standing of the concept. In Calleja's hands incorporation is a psychological intensification of involvement, set on ``their experientialist ontology'' after Lakoff and Johnson and on ``a view of consciousness as an internally generated construct'' after Dennett and Damasio (168--169). The reading I develop lifts it from that ground onto another: incorporation is the player's \textit{dwelling} in the designed \textit{Umwelt}---what Part I, after Heidegger, names ``the characteristic absorption of concern in its equipmental world'' (\textit{BT} 405)---the \textit{circumspective concern} in which the tool withdraws into its use, the hammer disappearing into the hammering itself (\textit{BT} 98). World-disclosedness simply \textit{is} this absorption of concern, redescribed for the virtual world. I must say plainly that this re-grounding is mine and not Calleja's: Heidegger is absent from his chapter on incorporation, and he reaches instead for cognitive science. Yet his own text invites the move twice over. He grants that ``incorporation is ultimately a metaphor like presence or immersion'' (3)---and a metaphor leaves its ontological debt unpaid; and there is a gap, plain on the page, between the experientialist ground he claims and the phenomenological words he in fact reaches for, the ``sense of habitation,'' the ``inhabiting a place.'' The re-grounding does not import a foreign vocabulary so much as pay for the one he had already spent.

What is \textit{elevated} is incorporation's reach: no longer the whole of the union but one of its two poles, joined to co-disclosedness within a single event. And the joining is prefigured in Calleja's own list---for one of the six dimensions, the shared, already names involvement with one's fellows, and so points past the world toward co-disclosedness. Here Calleja earns his place as an instrument and not an ornament, for his own writing yields a test of incorporation against mere involvement, which the method numbers and extends as the criteria R1--R7 and makes the gauge of world-disclosedness.\footnote{The incorporation criteria (Calleja, \textit{In-Game}, 169--173, 178; full articulation in the Calleja appendix): R1 assimilation; R2 embodiment; R3 simultaneity (the strictest gatekeeper, Calleja's own necessary condition); R4 the spatial-kinesthetic cornerstone; R5 blending; R6 temporal apex; R7 intrusion-integration. R3--R7 are Calleja's own incorporation-versus-involvement test; R1--R2 are drawn from the double axis.} The strictest gatekeeper is R3, simultaneity, which is Calleja's own necessary condition; the cornerstone is R4, the spatial-and-kinesthetic base ``without [which] incorporation cannot take place'' (170). The union-claim is thereby kept from mere assertion: it is read off conditions, and it returns a verdict---incorporation achieved, partial, merely-involved, or ruptured.

Why Calleja, and not another account of immersion or presence? The answer runs through his own correcting moves, not through a contest with rivals. He excises the magic circle, that supposed clean wall between the played world and the lived one---and the excision is exactly what a world real in kind demands, since what is real in kind cannot be sealed behind such a membrane. And he corrects the one-way metaphors of immersion and presence, which figure the player as merely lowered into a world; his double axis restores the two-way traffic---the world taken into awareness as much as the player given over into the world---that ``one in substrate, different in being'' requires.

\subsection{Consubstantiation and the Grammar of Motives: Burke at the Other Pole}

Incorporation joins the player to the world acted in. It does not join the player to those acted alongside; and the chain, which tells what moved the player, does not tell how that motion is afterward rendered, nor how its true source comes to be concealed. Two things are left unaccounted for, then---and neither is a \textit{residue} in the sense the term carries here, but each is a real gap: first, a consubstantial bond among several agents, a ``we'' that can itself \textit{source} an act, so that the deed springs from the bond and not from any single locus; second, a grammar by which a finished act is furnished with its motive, and by which that motive's source may be displaced. For both the dissertation turns to Kenneth Burke, and the turning is again a sublation.

What is \textit{preserved} is consubstantiation. ``In being identified with B,'' Burke writes, ``A is `substantially one' with a person other than himself. Yet at the same time he remains unique, an individual locus of motives. Thus he is both joined and separate'' (Burke, \textit{Rhetoric}, 21). That ``joined and separate'' is, word for word, the grammar of ``one in substrate, different in being,'' and its ground is act: ``substance\ldots{} was an act; and a way of life is an acting-together'' (21). The bond, moreover, is born of division and answers to it---``Identification is affirmed with earnestness precisely because there is division. Identification is compensatory to division'' (22)---and the player arrives divided in the most literal way, an actual body set over against the avatar acted as. What is \textit{transformed} is the engine of the bond. For Burke the ``we'' is worked by symbolic persuasion; in the method it is worked by shared desire---the ``we'' ratified at $A_3$, where the committed cognition of agents acting together fixes a common \textit{orekton}. What is \textit{elevated} is its place: consubstantiation becomes co-disclosedness, correlative to world-disclosedness, and folded into what Part I called the affective architecture of bodily-being-in-the-world-with-others.

The grammar of motives Burke supplies on his own warrant. He ``refers human actions to seven causes (\textit{aitia}): chance, nature, compulsion, habit, reason, anger, and desire'' (Burke, \textit{Grammar}, 292)---and that warrants reading the seven causes of \textit{Rhetorica} I.10 (1369a5--6) as a table of motive. The crosswalk from those seven to the pentad is mine, built on Burke's own move. His explicit map runs from the \textit{four} causes; I mark the difference rather than blur it. The pentad itself---Act, Scene, Agent, Agency, Purpose, with Attitude as the sixth term\footnote{The dramatistic pentad and its sixth term (Burke, \textit{Grammar}, xv; \textit{Rhetoric}, 42; full articulation in the Burke appendix): Act (``what was done''), Scene (``when or where it was done''), Agent (``who did it''), Agency (``how he did it''), Purpose (``why''), and Attitude, the sixth term, ``attitude being an incipient act.''}---the method takes as a grammar of motive-\textit{attribution}, ``a concern with the terms alone'' (xvi), not a psychology of the drives beneath an act: it names what a finished deed gets ascribed to, and lets one and the same deed be relocated from term to term. To it attaches a concealment diagnostic, which I frame narrowly and not as a blanket suspicion: ``in banishing the term, far from banishing its functions one merely conceals them,'' so that the analyst may ``detect its covert influence even in cases where it is overtly absent'' (21). The diagnostic catches a displaced motive---a source smuggled in under another term's name---and claims nothing more sweeping than that.

One last piece of evidence binds the two sources into one. Calleja, grounding incorporation in cognitive science, and Burke, grounding consubstantiation in symbolic action, arrive---each blind to the other, neither citing the other---at the very same figure of the joined-yet-separate: Calleja's two-way axis and Burke's ``both joined and separate.'' Two inquiries that set out from incompatible grounds and come to rest on one structure are not thereby proven right, but they are strongly corroborated; where the same form is deposited by hands that never met, it is likeliest that the form lies in the matter and not in the eye. That one structure is what the method names rhetorical incorporation---and world-disclosedness and co-disclosedness are not two lenses laid over a scene, but the two poles a single convergence brings to light.

\bigskip
\noindent\textit{Claims not grounded in the pack: none; every quotation is from a pack source. Three constructions are the dissertation's own and declared as such: the grounding of the world-/co-disclosedness union in the actualization of desire, consummated at $A_3 \rightarrow A_4$; the re-grounding of Calleja's incorporation and Burke's consubstantiation onto Heidegger's disclosedness pair; and the seven-causes\,$\rightarrow$\,pentad crosswalk.}

\end{document}
